% English translation of Jean-Luc Sauvageot, J. Operator Theory 9 (1983), 237--252.
% This file is self-contained: no external bibliography, body parts, or figures.
% Compile with pdfLaTeX (at least twice for all internal links).
% The requested filename contains TWO spaces before _GPT. A portable command is:
% pdflatex -jobname=translation "\input{1983 Sauvogeot_On the Relative Tensor Product of Hilbert Spaces\space\space_GPT.tex}"
% Rename translation.pdf to match this source filename if using that command.
\documentclass[fontsize=13pt,twoside,headings=small,bookmarkpackage=false]{scrartcl}
\usepackage[T1]{fontenc}
\usepackage[utf8]{inputenc}
\usepackage[english]{babel}
\usepackage{newpxtext}
\usepackage{amsmath,amsthm}
\usepackage{newpxmath}
\usepackage{geometry}
\geometry{paperwidth=176mm,paperheight=250mm,left=7mm,right=7mm,
  top=9.1mm,bottom=9.1mm,includehead,headheight=16pt,headsep=8pt,footskip=6pt}
\usepackage{fancyhdr}
\usepackage{tikz-cd}
\usepackage{enumitem}
\usepackage{microtype}
\usepackage{xcolor}
\definecolor{LinkColor}{HTML}{006A80}
\usepackage[unicode,colorlinks=true,linkcolor=LinkColor,citecolor=LinkColor,
  urlcolor=LinkColor,filecolor=LinkColor,linktoc=all,bookmarksnumbered=true,
  pdftitle={On the Relative Tensor Product of Hilbert Spaces},
  pdfauthor={Jean-Luc Sauvageot},pdfsubject={English translation of the 1983 article}]{hyperref}
\linespread{1.1}
\allowdisplaybreaks[4]
% Keep the opening expression with the first equality of long derivations;
% later align rows remain breakable, so display breaks actually work.
\setlength{\emergencystretch}{1.5em}
\setlength{\parindent}{1.25em}
\setlength{\parskip}{0pt}
\clubpenalty=10000
\widowpenalty=10000
\raggedbottom
\pagestyle{fancy}
\fancyhf{}
\fancyhead[LE,RO]{\normalfont\thepage}
\renewcommand{\headrulewidth}{0pt}
\renewcommand{\footrulewidth}{0pt}
\setkomafont{disposition}{\normalfont\bfseries}
\setkomafont{section}{\normalsize\bfseries}
\RedeclareSectionCommand[beforeskip=1.3\baselineskip,afterskip=.65\baselineskip]{section}
\setlist{nosep,leftmargin=1.8em}
\newcommand{\Lop}{\mathcal L}
\newcommand{\HH}{\mathcal H}
\newcommand{\ip}[2]{\langle #1,#2\rangle}
\newcounter{translatornote}
\newcommand{\tfnote}[1]{\stepcounter{translatornote}\begingroup
  \renewcommand{\thefootnote}{T\arabic{translatornote}}\footnote{#1}\endgroup}
\makeatletter
% Raise inline destinations so a jump leaves the full formula or diagram visible.
\newcommand{\raisedtarget}[1]{\Hy@raisedlink{\hyper@anchorstart{#1}\hyper@anchorend}}
\newcommand{\itemhead}[2]{\par\addvspace{.65\baselineskip}\noindent
  \phantomsection\def\@currentlabel{#1}\label{p:#1}\textbf{#1.\if\relax\detokenize{#2}\relax\else\ #2\fi}\enspace\ignorespaces}
\makeatother
\renewcommand{\proofname}{Proof}
% The source has no numbered displays. Preserve that choice; named internal
% targets are supplied without adding visible equation numbers.
\newcommand{\eqanchor}[1]{\hypertarget{eq:#1}{}}
\begin{document}
\setcounter{page}{1}
\thispagestyle{empty}
\begin{center}
{\small\textsc{J. Operator Theory} \textbf{9} (1983), 237--252\par}
\vspace{.8\baselineskip}
{\fontsize{17}{20}\selectfont\bfseries
Sur le produit tensoriel relatif d'espaces de Hilbert\par}
\vspace{.55\baselineskip}
{\fontsize{18}{21}\selectfont\bfseries
On the Relative Tensor Product of Hilbert Spaces\par}
\vspace{.7\baselineskip}
{\scshape Jean-Luc Sauvageot\par}
\vspace{.3\baselineskip}
{\footnotesize\copyright\ INCREST, 1983\par}
\end{center}

\section*{Introduction}
\addcontentsline{toc}{section}{Introduction}
\label{sec:intro}
The relative tensor product of Hilbert spaces is a procedure that assigns to a pair $(H,K)$ of Hilbert spaces, equipped respectively with a right $N$-module structure and a left $N$-module structure ($N$ being a fixed von Neumann algebra), a Hilbert space $H\otimes_N K$; it satisfies the following two requirements:
\begin{enumerate}[label=\arabic*)]
\item It is functorial; that is, if $(H',K')$ is another pair of Hilbert spaces with the same properties, one can form the tensor product $T_1\otimes_N T_2$ of two intertwining operators, $T_1\in\Lop_{N^\circ}(H,H')$, $T_2\in\Lop_N(K,K')$.
\item It ``cancels'' the $N$-$N$ bimodule $L^2(N)$ of the standard representation of $N$, in the sense that the functor $K\mapsto L^2(N)\otimes_N K$ (from the category of Hilbert spaces equipped with a normal, nondegenerate representation of $N$ to itself) is equivalent to the identity functor.
\end{enumerate}
It is easy to verify that such a procedure exists and is unique up to functorial equivalence. But this abstract approach is hardly usable, and the author (\cite{8}), in the case where $N$ is commutative, and then A.~Connes (\cite{2}), in full generality, give an explicit construction, using an auxiliary normal semifinite faithful (n.s.f.) weight $\nu$ on $N$, of the relative tensor product $H\otimes_\nu K$. This is a Hilbert space generated by elementary tensors $\xi\otimes_\nu\eta$, $\xi\in H$, $\eta\in K$, with $\xi$ or $\eta$ $\nu$-bounded. This provides computational machinery that allows a number of classical results in the theory of von Neumann algebras to be rewritten in terms of ``relative amplification.''

For example, the fundamental theorem of representation theory:
\begin{quote}
``Two faithful representations of a von Neumann algebra can be obtained from one another by an amplification followed by a reduction''
\end{quote}
(\cite{3}, I.4, Theorem~3) becomes
\begin{quote}
``Two faithful representations of a von Neumann algebra are obtained \emph{canonically} from one another by a relative amplification''
\end{quote}
(\hyperref[p:3.4]{Corollary~3.4} below).

Or again, the classical Hilbert--Schmidt theorem on the standard representation of type~I factors admits an elegant generalization:

The standard representation of a von Neumann algebra $M$ (given spatially on a Hilbert space $\HH$) is the relative amplification of $M$ on the Hilbert space $\HH\otimes_\nu\overline{\HH}$, the tensor product being taken over an algebra $N$ anti-isomorphic to the commutant of $M$ in $\Lop(\HH)$; the canonical involution is then the symmetry $J_\nu$ that assigns to the elementary tensor $\xi\otimes_\nu\overline\eta$ the tensor $\eta\otimes_\nu\overline\xi$; the positive cone $P^+$ is the closed cone generated by the tensors $\xi\otimes_\nu\overline\xi$, with $\xi$ $\nu^\circ$-bounded (\hyperref[p:3.1]{Proposition~3.1}; cf.\ \cite{8}, \S~I.3, for the case of discrete algebras).

Although it is now customary (cf.\ \cite{1}, \cite{7}, etc.) to use, mainly to economize on symbols, the terminology and notation of modules for the representation theory of a von Neumann algebra, we have preferred to fix, in a preliminary section (\hyperref[sec:0]{\S~0}), the terminology and notation used in the body of the article.

\hyperref[sec:1]{Section~1} is devoted to further results on $N$-modules, dealing mainly with the various ``operator-valued inner products'' underlying this type of structure.

\hyperref[sec:2]{Section~2} recalls the definition of the relative tensor product and proves its invariance under a change of the auxiliary weight $\nu$.

\hyperref[sec:3]{Section~3} contains the description of the standard representation announced above (the generalized Hilbert--Schmidt theorem), and its application to the study of functors from the representation theory of a von Neumann algebra $N$ to that of a von Neumann algebra $M$: we obtain a reformulation of M.~Rieffel's results (\cite{7}) in the simplified setting of ``Hilbert space'' bimodules.

\setcounter{section}{-1}
\section{Terminology and Notation}\label{sec:0}
\itemhead{0.1}{}
We fix a von Neumann algebra $N$ and denote its opposite algebra by $N^\circ$.

\itemhead{0.1.1}{}
A \emph{left $N$-module} $K$ is a Hilbert space equipped with a normal nondegenerate representation $\pi_N^K$ of $N$ in $\Lop(K)$. When there is no risk of confusion, we write $y\eta$ instead of $\pi_N^K(y)\eta$ for the result of the action of an element $y$ of $N$ on a vector $\eta$ of $K$.

\itemhead{0.1.2}{}
A \emph{right $N$-module} $H$ is a left $N^\circ$-module. If $y$ is an element of $N$ and $y^\circ$ the corresponding element of $N^\circ$, we write either $y^\circ\xi$ or $\xi y$ for the (right) action of $y$ on the vector $\xi$ of $H$.

\itemhead{0.1.3}{}
To the left $N$-module $K$ there corresponds canonically the right $N$-module $\overline K$, defined on the conjugate space of $K$ by
\eqanchor{conjugate-action}
\[
\overline\eta y=\overline{y^*\eta},\qquad y\in N,\quad\eta\in K.
\]
We also denote by $H\mapsto\overline H$ the inverse correspondence between right $N$-modules and left $N$-modules.

\itemhead{0.1.4}{}
If $K_1$ and $K_2$ are two left $N$-modules, we set
\[
\Lop_N(K_1,K_2)=\{T\in\Lop(K_1,K_2)\mid Ty\eta=yT\eta,\ \forall\eta\in K_1,\ \forall y\in N\};
\]
we write $\Lop_N(K_1)$ instead of $\Lop_N(K_1,K_1)$.
If $H_1$ and $H_2$ are two right $N$-modules, we have
\begin{align*}
\Lop_{N^\circ}(H_1,H_2)=\{T\in\Lop(H_1,H_2)\mid{}& (T\xi)y=T(\xi y),\\
&\forall\xi\in H_1,\ \forall y\in N\}.
\end{align*}

\itemhead{0.2}{}
Let $M$ be a second von Neumann algebra.

\itemhead{0.2.1}{}
An \emph{$M$-$N$ bimodule} is a Hilbert space $\HH$ equipped with a left $M$-module structure and a right $N$-module structure that commute, that is, such that
\[
(x\xi)y=x(\xi y),\qquad\forall\xi\in\HH,\quad\forall x\in M,\quad\forall y\in N.
\]

\itemhead{0.2.2}{}
By definition, an $M$-$N$ bimodule is also an $N^\circ$-$M^\circ$ bimodule, and we have
\[
y^\circ\xi x^\circ=x\xi y,\qquad\forall\xi\in\HH,\quad\forall x\in M,\quad\forall y\in N.
\]

\itemhead{0.2.3}{}
If $\HH$ is an $M$-$N$ bimodule, then by \hyperref[p:0.1.3]{0.1.3} the conjugate Hilbert space $\overline\HH$ is an $N$-$M$ bimodule, and we have
\[
y\overline\xi x=\overline{x^*\xi y^*},\qquad\forall\xi\in\HH,\quad\forall x\in M,\quad\forall y\in N.
\]

\itemhead{0.3.1}{}
For every von Neumann algebra $N$ we fix a standard $N$-$N$ bimodule: this is the Hilbert space $L^2(N)$ (unique up to isomorphism), equipped with the natural representation of $N$, the antilinear involution $J$, and the self-dual cone $L^2(N)_+$ satisfying the axioms of \cite{4}, Definition~2.1. The right $N$-module structure is of course given by
\[
\xi y=Jy^*J\xi,\qquad\xi\in L^2(N),\quad y\in N.
\]

\itemhead{0.3.2}{}
If $\nu$ is an n.s.f.\ weight on $N$, there is a unique isometric $N$-module morphism $\Lambda_\nu$ from the left ideal $\mathcal N_\nu$ into $L^2(N)$ satisfying
\[
yJ\Lambda_\nu(y)\in L^2(N)_+,\qquad\forall y\in\mathcal N_\nu
\]
(by ``isometric $N$-module morphism'' we mean the following properties:
\begin{align*}
\Lambda_\nu(yy')&=y\Lambda_\nu(y'),\qquad\forall y\in N,\quad\forall y'\in\mathcal N_\nu,\\
\|\Lambda_\nu(y)\|^2&=\nu(y^*y),\qquad\forall y\in\mathcal N_\nu).
\end{align*}

\itemhead{0.3.3}{}
For $L^2(N^\circ)$ we choose the $N^\circ$-$N^\circ$ bimodule $L^2(N)$, with the same involution $J$ and the same self-dual cone $L^2(N)_+=L^2(N^\circ)_+$. With this choice, if $\nu^\circ$ is the weight on $N^\circ$ corresponding to the n.s.f.\ weight $\nu$ on $N$, it is easy to verify that
\[
\Lambda_{\nu^\circ}(y^\circ)=J\Lambda_\nu(y^*),\qquad
\forall y\in\mathcal N_\nu^*=(\mathcal N_{\nu^\circ})^\circ.
\]

\itemhead{0.4.1}{}
Let $\nu$ be an n.s.f.\ weight on $N$. Recall (cf.\ \cite{1}) that $D(K,\nu)$ denotes the set of $\nu$-bounded vectors of a left $N$-module $K$, that is, vectors $\eta$ of $K$ for which there exists $R^\nu(\eta)$ in $\Lop_N(L^2(N),K)$ satisfying
\eqanchor{bounded-vector}
\[
R^\nu(\eta)\Lambda_\nu(y)=y\eta,\qquad\forall y\in\mathcal N_\nu.
\]
To a pair $(\eta,\eta')$ of $\nu$-bounded vectors in $K$ is associated the operator $\theta^\nu(\eta,\eta')=R^\nu(\eta)R^\nu(\eta')^*$ in $\Lop_N(K)$.

\itemhead{0.4.2}{}
If $H$ is a right $N$-module and $\xi$ in $H$ is $\nu^\circ$-bounded, $R^{\nu^\circ}(\xi)$ is the operator in $\Lop_{N^\circ}(L^2(N),H)$ characterized by
\[
R^{\nu^\circ}(\xi)J\Lambda_\nu(y)=\xi y^*,\qquad\forall y\in\mathcal N_\nu.
\]

\itemhead{0.5}{}
If $\nu$ is an n.s.f.\ weight on $N$ and $\lambda$ a complex number, the \emph{domain} of $\sigma_\lambda^\nu$, denoted by $\mathcal D(\sigma_\lambda^\nu)$, is the set of $y\in N$ such that the map $\mathbb R\ni t\mapsto\sigma_t^\nu(y)$ admits a continuous extension to the strip $0\leq\operatorname{Im}z\leq\operatorname{Im}\lambda$ (or $\operatorname{Im}\lambda\leq\operatorname{Im}z\leq0$), analytic in its interior. Then $\sigma_\lambda^\nu(y)$ is the value of this extension at $z=\lambda$ and is characterized by
\eqanchor{analytic-modular-action}
\[
\sigma_\lambda^\nu(y)\Delta_\nu^{i\lambda}\subset\Delta_\nu^{i\lambda}y.
\]


\section{Further Results on \texorpdfstring{$N$}{N}-Modules}\label{sec:1}

We fix a von Neumann algebra $N$ and a normal semifinite faithful weight $\nu$ on $N$; $\nu^\circ$ is the corresponding weight on the algebra $N^\circ$; $\mathcal U_0$ denotes the Tomita algebra associated with the weight $\nu$.

Throughout this section, $H$ denotes a right $N$-module and $K$ a left $N$-module.

\itemhead{1.1}{}
We equip the vector space $D(H,\nu^\circ)$ with an ``$N$-valued inner product'' as follows:

If $\xi_1$ and $\xi_2$ are two $\nu^\circ$-bounded vectors in $H$, then $R^{\nu^\circ}(\xi_2)^*R^{\nu^\circ}(\xi_1)$ is an endomorphism of $L^2(N)$ which commutes with the right action of $N$, and consequently defines an element $\ip{\xi_1}{\xi_2}_{\nu^\circ}$ of $N$ characterized by
\[
 \ip{\xi_1}{\xi_2}_{\nu^\circ}\zeta
 =R^{\nu^\circ}(\xi_2)^*R^{\nu^\circ}(\xi_1)\zeta,
 \qquad \forall\zeta\in L^2(N).
\]
We also equip $D(K,\nu)$ with an $N$-valued inner product by setting, for $\nu$-bounded $\eta_1,\eta_2$,
\[
 \ip{\eta_1}{\eta_2}_\nu
 =\ip{\overline{\eta_2}}{\overline{\eta_1}}_{\nu^\circ},
\]
that is,
\[
 \ip{\eta_1}{\eta_2}_\nu\zeta
 =JR^\nu(\eta_1)^*R^\nu(\eta_2)J\zeta,
 \qquad \forall\zeta\in L^2(N).
\]

\itemhead{1.2}{}
If $\nu$ is a trace, it is easy to check that a vector $\eta$ in $K$ is $\nu$-bounded if and only if $\frac{d\omega_{\eta,\eta}}{d\nu}$ is a bounded operator, and that we have\tfnote{The original writes $d\tau$ in both denominators in this paragraph although the trace is denoted by $\nu$. The notation has been made consistent.}
\[
 \ip{\eta_1}{\eta_2}_\nu
 =\frac{d\omega_{\eta_1,\eta_2}}{d\nu}
\]
(cf.~\cite{8}, Lemma~I.1; $\omega_{\eta_1,\eta_2}$ is the linear functional $N\ni y\mapsto\ip{y\eta_1}{\eta_2}$).

Then $D(K,\nu)$ is a ``pre-Hilbert $N$-module'' in the sense of~\cite{6}.

\itemhead{1.3}{}
In the general case, $D(K,\nu)$ is invariant under the action of $\mathcal D(\sigma^\nu_{i/2})$, and we have
\begin{align*}
 \ip{y\eta_1}{\eta_2}_\nu
 &=\sigma^\nu_{i/2}(y)\ip{\eta_1}{\eta_2}_\nu,\\
 &\qquad\forall\eta_1,\eta_2\in D(K,\nu),\quad
 \forall y\in\mathcal D(\sigma^\nu_{i/2});
\end{align*}
likewise, $D(H,\nu^\circ)$ is invariant under $\mathcal D(\sigma^\nu_{-i/2})$, and we have\tfnote{In the second formula of the original, the quantified vectors are printed as $\eta_1,\eta_2$; they have been corrected to $\xi_1,\xi_2$.}
\begin{align*}
 \ip{\xi_1y}{\xi_2}_{\nu^\circ}
 &=\ip{\xi_1}{\xi_2}_{\nu^\circ}\sigma^\nu_{-i/2}(y),\\
 &\qquad\forall\xi_1,\xi_2\in D(H,\nu^\circ),\quad
 \forall y\in\mathcal D(\sigma^\nu_{-i/2}).
\end{align*}

\itemhead{1.4}{}
If $H$ is the right $N$-module $L^2(N)$, then $D(H,\nu^\circ)$ is identified with $\Lambda_\nu(\mathcal N_\nu)$, and we have
\[
 \ip{\Lambda_\nu(y_1)}{\Lambda_\nu(y_2)}_{\nu^\circ}
 =y_2^*y_1,\qquad\forall y_1,y_2\in\mathcal N_\nu.
\]
If $K$ is the left $N$-module $L^2(N)$, then $D(K,\nu)$ is identified with $J\Lambda_\nu(\mathcal N_\nu)$, and we have
\[
 \ip{J\Lambda_\nu(y'_1)}{J\Lambda_\nu(y'_2)}_\nu
 ={y'_1}^*y'_2,\qquad\forall y'_1,y'_2\in\mathcal N_\nu.
\]

\itemhead{1.5}{Lemma.}
\textit{Let $\xi_1$ and $\xi_2$ be two $\nu^\circ$-bounded vectors in $H$, and $\eta_1$ and $\eta_2$ two $\nu$-bounded vectors in $K$. Then:}
\begin{align*}
 \text{a)}\quad&\ip{\xi_1}{\xi_2}_{\nu^\circ}\in\mathcal N_\nu,
 \quad \Lambda_\nu(\ip{\xi_1}{\xi_2}_{\nu^\circ})
 =R^{\nu^\circ}(\xi_2)^*\xi_1;\\
 \text{b)}\quad&\raisedtarget{eq:1.5b}\ip{\eta_1}{\eta_2}_\nu\in\mathcal N_\nu,
 \quad \Lambda_\nu(\ip{\eta_1}{\eta_2}_\nu)
 =JR^\nu(\eta_1)^*\eta_2;\\
 \text{c)}\quad&\raisedtarget{eq:1.5c}\ip{\xi_1\ip{\eta_1}{\eta_2}_\nu}{\xi_2}
 =\ip{\ip{\xi_1}{\xi_2}_{\nu^\circ}\eta_1}{\eta_2}.
\end{align*}

\begin{proof}
\hypertarget{proof:1.5a}{}a) By~\cite{1}, Lemma~4, $\ip{\xi_1}{\xi_2}_{\nu^\circ}\in\mathcal M_\nu\subset\mathcal N_\nu$, and we have
\[
 \nu(\ip{\xi_1}{\xi_2}_{\nu^\circ})=\ip{\xi_1}{\xi_2}.
\]
For $y\in\mathcal U_0$, we compute
\begin{align*}
 &\ip{\Lambda_\nu(\ip{\xi_1}{\xi_2}_{\nu^\circ})}{J\Lambda_\nu(y)}\\*
 &\quad=\ip{\Delta_\nu^{1/2}\Lambda_\nu(y)}{\Lambda_\nu(\ip{\xi_2}{\xi_1}_{\nu^\circ})}\\
 &\quad=\ip{\Lambda_\nu(\sigma^\nu_{-i/2}(y))}{\Lambda_\nu(\ip{\xi_2}{\xi_1}_{\nu^\circ})}\\
 &\quad=\nu\bigl(\ip{\xi_1}{\xi_2}_{\nu^\circ}\sigma^\nu_{-i/2}(y)\bigr)\\
 &\quad=\ip{\xi_1y}{\xi_2}
 \qquad\text{(by \hyperref[p:1.3]{1.3})}\\
 &\quad=\ip{\xi_1}{R^{\nu^\circ}(\xi_2)J\Lambda_\nu(y)}\\
 &\quad=\ip{R^{\nu^\circ}(\xi_2)^*\xi_1}{J\Lambda_\nu(y)}.
\end{align*}
b) This is proved in the same way as \hyperlink{proof:1.5a}{a)}.

Assertion \hyperlink{eq:1.5c}{c)} is stated in~\cite{2}, Proposition~12~b). To prove it, we compute
\begin{align*}
 \ip{\xi_1\ip{\eta_1}{\eta_2}_\nu}{\xi_2}
 &=\ip{J\Lambda_\nu(\ip{\eta_2}{\eta_1}_\nu)}{R^{\nu^\circ}(\xi_1)^*\xi_2}\\
 &=\ip{R^\nu(\eta_2)^*\eta_1}{\Lambda_\nu(\ip{\xi_2}{\xi_1}_{\nu^\circ})}\\
 &=\ip{\eta_1}{\ip{\xi_2}{\xi_1}_{\nu^\circ}\eta_2}.
\end{align*}
\end{proof}

\itemhead{1.6}{Lemma.}
\textit{Let $\psi$ be a normal semifinite faithful weight on $\Lop_N(K)$ and $\lambda$ a complex number.}

\textit{a) For $y\in\mathcal D(\sigma^\nu_{-i\lambda})$, we have the inclusion}
\eqanchor{1.6a}
\[
 \frac{d\nu^\lambda}{d\psi}\,y
 \supset\sigma^\nu_{-i\lambda}(y)\frac{d\nu^\lambda}{d\psi}.
\]
\textit{\raisedtarget{part:1.6b}b) If $\eta\in\mathcal D\bigl(\frac{d\nu^\lambda}{d\psi}\bigr)$ is $\nu$-bounded and $\frac{d\nu^\lambda}{d\psi}\eta$ is also $\nu$-bounded, we have the inclusion}
\[
 \frac{d\nu^\lambda}{d\psi}R^\nu(\eta)
 \supset R^\nu\left(\frac{d\nu^\lambda}{d\psi}\eta\right)\Delta_\nu^\lambda.
\]
\textit{c) If $\eta$ and $\eta'$ are as in \hyperlink{part:1.6b}{b)}, then $\bigl\langle\frac{d\nu^\lambda}{d\psi}\eta,\eta'\bigr\rangle_\nu$ belongs to the domain of $\sigma^\nu_{i\lambda}$, and we have}
\[
 \sigma^\nu_{i\lambda}\left(\left\langle
 \frac{d\nu^\lambda}{d\psi}\eta,\eta'\right\rangle_\nu\right)
 =\left\langle\eta,\frac{d\nu^{\bar\lambda}}{d\psi}\eta'\right\rangle_\nu.
\]

\begin{proof}
a) Let $\eta,\eta'\in\mathcal D\bigl(\frac{d\nu^\lambda}{d\psi}\bigr)$. The two functions of the complex variable
\[
 z\longmapsto\left\langle y\eta,
 \frac{d\nu^{\bar z}}{d\psi}\eta'\right\rangle
\]
and
\[
 z\longmapsto\left\langle\frac{d\nu^z}{d\psi}\eta,
 \sigma^\nu_{i\bar z}(y^*)\eta'\right\rangle
\]
are continuous in the strip $0\leq\operatorname{Re}z\leq\operatorname{Re}\lambda$ (or $\operatorname{Re}\lambda\leq\operatorname{Re}z\leq0$), analytic in its interior, and coincide for $z=it$, $t\in\mathbb R$. They are therefore equal everywhere, in particular for $z=\lambda$, which proves that
\[
 y\eta\in\mathcal D\left[\left(\frac{d\nu^{\bar\lambda}}{d\psi}\right)^*\right]
 =\mathcal D\left(\frac{d\nu^\lambda}{d\psi}\right),
\]
and gives the required equality.

b) It suffices to verify the inclusion on $\Lambda_\nu(\mathcal U_0)$, which is a core for $\Delta_\nu^\lambda$; the assertion then follows from \hyperlink{eq:1.6a}{a)}.

c) Applying \hyperlink{part:1.6b}{b)} twice gives the following equalities and inclusions between operators on $L^2(N)$:
\begin{align*}
 &\left\langle\eta,\frac{d\nu^{\bar\lambda}}{d\psi}\eta'\right\rangle_\nu\Delta_\nu^{-\lambda}\\*
 &\quad=JR^\nu(\eta)^*R^\nu\left(\frac{d\nu^{\bar\lambda}}{d\psi}\eta'\right)J\Delta_\nu^{-\lambda}\\
 &\quad=JR^\nu(\eta)^*R^\nu\left(\frac{d\nu^{\bar\lambda}}{d\psi}\eta'\right)\Delta_\nu^{\bar\lambda}J\\
 &\quad\subset JR^\nu(\eta)^*\frac{d\nu^{\bar\lambda}}{d\psi}R^\nu(\eta')J\\
 &\quad\subset J\Delta_\nu^{\bar\lambda}R^\nu\left(\frac{d\nu^\lambda}{d\psi}\eta\right)^*R^\nu(\eta')J\\
 &\quad=\Delta_\nu^{-\lambda}\left\langle\frac{d\nu^\lambda}{d\psi}\eta,\eta'\right\rangle_\nu.
\end{align*}
The result follows.
\end{proof}

\itemhead{1.7}{Lemma.}
\textit{Let $H$ be a right $N$-module, $\varphi$ a normal semifinite faithful weight on $\Lop_{N^\circ}(H)$, and $\xi$ a vector in the domain of $\bigl(\frac{d\nu^\circ}{d\varphi}\bigr)^{1/2}$.}

\textit{Then $\xi$ is $\varphi$-bounded if and only if $\bigl(\frac{d\nu^\circ}{d\varphi}\bigr)^{1/2}\xi$ is $\nu^\circ$-bounded; and we have the following equality in $N$:}\tfnote{The original omits $\xi$ after the operator in this boundedness condition. The opposite-algebra superscript on $\theta^\varphi$ is also restored in the second calculation in the proof.}
\[
 \theta^\varphi(\xi,\xi)^\circ
 =\left\langle\left(\frac{d\nu^\circ}{d\varphi}\right)^{1/2}\xi,
 \left(\frac{d\nu^\circ}{d\varphi}\right)^{1/2}\xi\right\rangle_{\nu^\circ}.
\]

\begin{proof}
Suppose that $\xi$ is $\varphi$-bounded. For $y\in\mathcal U_0$, we have
\begin{align*}
 \nu\bigl(y^*\theta^\varphi(\xi,\xi)^\circ y\bigr)
 &=\nu^\circ\bigl(\theta^\varphi(\xi y,\xi y)\bigr)\\
 &=\left\|\left(\frac{d\nu^\circ}{d\varphi}\right)^{1/2}(\xi y)\right\|^2\\
 &=\left\|\left(\left(\frac{d\nu^\circ}{d\varphi}\right)^{1/2}\xi\right)
 \sigma^\nu_{i/2}(y)\right\|^2,
 \quad\text{by \hyperlink{eq:1.6a}{1.6~a)}};
\end{align*}
that is,
\[
 \forall y\in\mathcal U_0,\qquad
 \left\|\left(\left(\frac{d\nu^\circ}{d\varphi}\right)^{1/2}\xi\right)y\right\|^2
 \leq\|\theta^\varphi(\xi,\xi)\|\,\nu(yy^*),
\]
which ensures that $\bigl(\frac{d\nu^\circ}{d\varphi}\bigr)^{1/2}\xi$ is $\nu^\circ$-bounded.

We then have
\begin{align*}
 &\nu\bigl(y^*\theta^\varphi(\xi,\xi)^\circ y\bigr)\\*
 &\quad=\nu\left(\left\langle
 \left(\left(\frac{d\nu^\circ}{d\varphi}\right)^{1/2}\xi\right)\sigma^\nu_{i/2}(y),
 \left(\left(\frac{d\nu^\circ}{d\varphi}\right)^{1/2}\xi\right)\sigma^\nu_{i/2}(y)
 \right\rangle_{\nu^\circ}\right)\\
 &\quad=\nu\left(y^*\left\langle
 \left(\frac{d\nu^\circ}{d\varphi}\right)^{1/2}\xi,
 \left(\frac{d\nu^\circ}{d\varphi}\right)^{1/2}\xi
 \right\rangle_{\nu^\circ}y\right),
 \quad\text{by \hyperref[p:1.3]{1.3}}.
\end{align*}
This gives the required equality.

Conversely, if $\bigl(\frac{d\nu^\circ}{d\varphi}\bigr)^{1/2}\xi$ is assumed to be $\nu^\circ$-bounded, the preceding computation ensures that
\[
 \xi=\left(\frac{d\varphi}{d\nu^\circ}\right)^{1/2}
 \left(\left(\frac{d\nu^\circ}{d\varphi}\right)^{1/2}\xi\right)
\]
is $\varphi$-bounded.
\end{proof}

\itemhead{1.8}{Remark.}
To conclude, consider an $M$--$N$ bimodule $\HH$ such that we have exactly $L_M(\HH)=L_{N^\circ}(\HH)'$.

If $\varphi$ is a normal semifinite faithful weight on $M$, we set\tfnote{The arguments of $D_\infty$ are printed inconsistently in the original, first with $\psi$ in place of $\varphi$ and then in reversed order. They are written throughout as $(\HH,\varphi,\nu)$ here. The subscript $\infty$, omitted in the final completion statement, is also restored.}
\begin{align*}
 D_\infty(\HH,\varphi,\nu)
 &=\bigcap_{\lambda\in\mathbb R}
 \mathcal D\left(\left(\frac{d\nu^\circ}{d\varphi}\right)^\lambda\right)
 \cap D(\HH,\nu^\circ),\\
 M_\infty(\varphi)&=\bigcap_{\lambda\in\mathbb R}\mathcal D(\sigma^\varphi_{i\lambda}),\\
 N_\infty(\nu)&=\bigcap_{\lambda\in\mathbb R}\mathcal D(\sigma^\nu_{i\lambda}).
\end{align*}
First, by~\hyperref[p:1.7]{1.7}, we have
\[
 D_\infty(\HH,\varphi,\nu)
 =\bigcap_{\lambda\in\mathbb R}
 \mathcal D\left(\left(\frac{d\nu^\circ}{d\varphi}\right)^\lambda\right)
 \cap D(\HH,\varphi).
\]
Next, by~\hyperlink{eq:1.6a}{1.6~a)}, $D_\infty(\HH,\varphi,\nu)$ is an $M_\infty(\varphi)$--$N_\infty(\nu)$ bimodule, on which we now consider the inner products\tfnote{The subscript $\nu$ on the first inner product and the second equals sign in the second formula, both missing in the original, have been restored.}
\begin{align*}
 (\xi\mid\eta)_\nu
 &=\theta^\varphi(\xi,\eta)^\circ
 =\left\langle\left(\frac{d\nu^\circ}{d\varphi}\right)^{1/2}\xi,
 \left(\frac{d\nu^\circ}{d\varphi}\right)^{1/2}\eta\right\rangle_{\nu^\circ},\\
 (\xi\mid\eta)_\varphi
 &=\theta^{\nu^\circ}(\xi,\eta)
 =\left\langle\left(\frac{d\varphi}{d\nu^\circ}\right)^{1/2}\xi,
 \left(\frac{d\varphi}{d\nu^\circ}\right)^{1/2}\eta\right\rangle_\varphi,
\end{align*}
which take values in $N$ and $M$, respectively. (By~\cite{5}, II, Proposition~2, $D_\infty(\HH,\varphi,\nu)$ is dense in $\HH$.)

By completing $D_\infty(\HH,\varphi,\nu)$ for the family of seminorms
$\{\xi\mapsto\omega((\xi\mid\xi)_\nu)^{1/2};\ \omega\in N_*^+\}$,
we obtain an $M$--$N$ imprimitivity bimodule in the sense of~\cite{7}, which is easily checked to be the ``self-adjoint equivalence module'' giving the Morita equivalence between $M$ and $N$ in Theorem~8.15 of~\cite{7}.

The notion of the ``tensor product of $N$-modules'' makes it possible (for $W^*$-algebras only) to avoid recourse to imprimitivity bimodules, whose entire structure is, in a certain sense, already contained in the more classical and simpler bimodules considered here. In particular, these implicitly contain the ``operator-valued inner products'' used extensively in~\cite{1} and~\cite{7}.


\section{Tensor Product of \texorpdfstring{$N$}{N}-Modules}\label{sec:2}

\itemhead{2.1}{Definition.}
(\cite{2}; cf.~\cite{8} for the commutative case.) Let $N$ be a von Neumann algebra, $\nu$ a faithful normal semifinite weight on $N$, $H$ a right $N$-module, and $K$ a left $N$-module. The \emph{tensor product over $N$, modulo $\nu$} of $H$ and $K$, denoted by $H\otimes_\nu K$, is the Hilbert space obtained by taking the separated completion of the algebraic tensor product
\[
D(H,\nu^\circ)\odot K
\]
equipped with the inner product defined on elementary tensors by
\[
\ip{\xi_1\odot\eta_1}{\xi_2\odot\eta_2}
=\ip{\ip{\xi_1}{\xi_2}_{\nu^\circ}\eta_1}{\eta_2}.
\]
By \hyperlink{eq:1.5c}{1.5(c)}, $H\otimes_\nu K$ is also the Hilbert space obtained by taking the separated completion of $H\odot D(K,\nu)$ for the inner product
\[
\ip{\xi_1\odot\eta_1}{\xi_2\odot\eta_2}
=\ip{\xi_1\ip{\eta_1}{\eta_2}_\nu}{\xi_2}.
\]
We denote by $\xi\otimes_\nu\eta$ the canonical image in $H\otimes_\nu K$ of the tensor $\xi\odot\eta$, where $\xi\in D(H,\nu^\circ)$ and $\eta\in K$ (or $\xi\in H$ and $\eta\in D(K,\nu)$).

\itemhead{2.2}{Remark.}
(a) For fixed $\xi_0\in D(H,\nu^\circ)$ (respectively, fixed $\eta_0\in D(K,\nu)$), the linear map
\[
K\ni\eta\longmapsto\xi_0\otimes_\nu\eta\in H\otimes_\nu K
\]
(respectively,
\[
H\ni\xi\longmapsto\xi\otimes_\nu\eta_0\in H\otimes_\nu K)
\]
is continuous.

In particular, if $D_1$ is a dense vector subspace of $D(H,\nu^\circ)$ (respectively, of $H$), and $D_2$ is a dense vector subspace of $K$ (respectively, of $D(K,\nu)$), then $D_1\otimes_\nu D_2$ is a dense vector subspace of $H\otimes_\nu K$.

(b) By \cite[Remark~16(b)]{2}, or by \hyperref[p:1.3]{1.3} above, we have
\begin{align*}
\xi y\otimes_\nu\eta
&=\xi\otimes_\nu\sigma^\nu_{-i/2}(y)\eta,\\
&\hspace{-1em}\forall\xi\in H,\quad\forall\eta\in D(K,\nu),\quad
\forall y\in\mathcal D(\sigma^\nu_{-i/2}).
\end{align*}
The functoriality of this construction follows easily from \cite[Proposition~12]{2}:

\itemhead{2.3}{Lemma.}
Let $H$ and $H'$ be two right $N$-modules, $K$ and $K'$ two left $N$-modules, and $T_1$ and $T_2$ two operators in $\Lop_{N^\circ}(H,H')$ and $\Lop_N(K,K')$, respectively.

There is then a continuous operator $T_1\otimes_\nu T_2$ from $H\otimes_\nu K$ to $H'\otimes_\nu K'$ satisfying
\[
(T_1\otimes_\nu T_2)(\xi\otimes_\nu\eta)
=T_1\xi\otimes_\nu T_2\eta,
\qquad\forall\xi\in D(H,\nu^\circ),\quad\forall\eta\in K
\]
(or $\forall\xi\in H$, $\forall\eta\in D(K,\nu)$).

In particular, $H\otimes_\nu K$ is an $\Lop_{N^\circ}(H)$--$\Lop_N(K)^\circ$ bimodule.

\itemhead{2.4}{}
Lemma~15 of \cite{2} can be reformulated as follows:

Let $H$ be a right $N$-module and $K$ a left $N$-module.

(a) There is one and only one $\Lop_{N^\circ}(H)$--$N$ bimodule isomorphism $U_H^\nu$ from $H\otimes_\nu L^2(N)$ onto $H$, satisfying\tfnote{The source omits the opposite-algebra sign in $\Lop_{N^\circ}(H)$ here and in \hyperlink{part:2.6a}{Proposition~2.6(a)}. It is restored because $H$ is a right $N$-module.}
\[
U_H^\nu\bigl(\xi\otimes_\nu J\Lambda_\nu(y)\bigr)
=\xi y^*,\qquad\forall\xi\in H,\quad\forall y\in\mathcal N_\nu.
\]

(b) There is one and only one $N$--$\Lop_N(K)^\circ$ bimodule isomorphism $V_K^\nu$ from $L^2(N)\otimes_\nu K$ onto $K$, satisfying
\[
V_K^\nu\bigl(\Lambda_\nu(y)\otimes_\nu\eta\bigr)
=y\eta,\qquad\forall y\in\mathcal N_\nu,\quad\forall\eta\in K.
\]
Notice that, modulo the identifications of \hyperref[sec:0]{Section~0}, we have
\[
U_{L^2(N)}^\nu=V_{L^2(N)}^\nu.
\]

\itemhead{2.5}{Lemma.}
With the notation of the statements of \hyperref[p:2.3]{Lemma~2.3} and \hyperref[p:2.4]{2.4}, the following two diagrams commute:\tfnote{In the first diagram, the original lower-left object is printed as $H\otimes_\nu L^2(N)$. The prime on $H'$ has been restored, as required by the arrows $T_1\otimes_\nu i$ and $U_{H'}^\nu$.}
\[
\begin{tikzcd}[column sep=large,row sep=large]
H\otimes_\nu L^2(N) \arrow[r,"U_H^\nu"] \arrow[d,"T_1\otimes_\nu i"']
&H\arrow[d,"T_1"]\\
H'\otimes_\nu L^2(N)\arrow[r,"U_{H'}^\nu"']&H'
\end{tikzcd}
\qquad
\begin{tikzcd}[column sep=large,row sep=large]
L^2(N)\otimes_\nu K\arrow[r,"V_K^\nu"]\arrow[d,"i\otimes_\nu T_2"']
&K\arrow[d,"T_2"]\\
L^2(N)\otimes_\nu K'\arrow[r,"V_{K'}^\nu"']&K'
\end{tikzcd}
\]
($i$ denotes the identity of $L^2(N)$).

The verification of the lemma is trivial.

\itemhead{2.6}{Proposition.}
\hypertarget{part:2.6a}{}(a) Let $H$ be a right $N$-module, $K$ a left $N$-module, and $\nu$ and $\nu'$ two faithful normal semifinite weights on $N$.

There is a unique $\Lop_{N^\circ}(H)$--$\Lop_N(K)^\circ$ bimodule isomorphism $U_{H,K}^{\nu,\nu'}$ from $H\otimes_\nu K$ onto $H\otimes_{\nu'}K$ satisfying the following property:\tfnote{The original prose prints the target as $H\otimes_\nu K$. The prime on $\nu'$ has been restored, in agreement with the diagram immediately below.} for all $t_1\in\Lop_{N^\circ}(H,L^2(N))$ and all $t_2\in\Lop_N(K,L^2(N))$, the diagram
\[
\begin{tikzcd}[column sep=4em,row sep=large]
H\otimes_\nu K\arrow[rr,"U_{H,K}^{\nu,\nu'}"]\arrow[d,"t_1\otimes_\nu t_2"']
&&H\otimes_{\nu'}K\arrow[d,"t_1\otimes_{\nu'}t_2"]\\
L^2(N)\otimes_\nu L^2(N)\arrow[r,"U_{L^2(N)}^\nu"']
&L^2(N)
&L^2(N)\otimes_{\nu'}L^2(N)\arrow[l,"U_{L^2(N)}^{\nu'}"]
\end{tikzcd}
\]
commutes.

(b) If $H'$ and $K'$ are two $N$-modules, right and left respectively, and $T_1$ and $T_2$ are two operators in $\Lop_{N^\circ}(H,H')$ and $\Lop_N(K,K')$, respectively, then the diagram\raisedtarget{diagram:2.6b}
\[
\begin{tikzcd}[column sep=huge,row sep=large]
H\otimes_\nu K\arrow[r,"U_{H,K}^{\nu,\nu'}"]\arrow[d,"T_1\otimes_\nu T_2"']
&H\otimes_{\nu'}K\arrow[d,"T_1\otimes_{\nu'}T_2"]\\
H'\otimes_\nu K'\arrow[r,"U_{H',K'}^{\nu,\nu'}"']&H'\otimes_{\nu'}K'
\end{tikzcd}
\]
commutes.

\begin{proof}
(a) \emph{Existence.} Fix $H$ and suppose that $U_{H,K_0}^{\nu,\nu'}$ exists with the desired properties for a left $N$-module $K_0$. Then, by \hyperref[p:2.3]{Lemma~2.3}, $U_{H,K}^{\nu,\nu'}$ will exist if $K$ is an amplification of $K_0$, or a reduction of such an amplification by a projection commuting with the action of $N$.

For existence we may therefore reduce to the case $K=L^2(N)$; it then suffices to take
\[
U_{H,L^2(N)}^{\nu,\nu'}=(U_H^{\nu'})^*U_H^\nu
\]
and to verify, for $t_1\in\Lop_{N^\circ}(H,L^2(N))$ and $t_2\in\Lop_N(K,L^2(N))$, that the diagram
\[
\begin{tikzcd}[column sep=4em,row sep=large]
H\otimes_\nu L^2(N)\arrow[r,"U_H^\nu"]\arrow[d,"t_1\otimes_\nu t_2"']
&H\arrow[d,"t_1t_2"]
&H\otimes_{\nu'}L^2(N)\arrow[l,"U_H^{\nu'}"']\arrow[d,"t_1\otimes_{\nu'}t_2"]\\
L^2(N)\otimes_\nu L^2(N)\arrow[r,"U_{L^2(N)}^\nu"']
&L^2(N)
&L^2(N)\otimes_{\nu'}L^2(N)\arrow[l,"U_{L^2(N)}^{\nu'}"]
\end{tikzcd}
\]
commutes.

\emph{Uniqueness.} Let $\xi\in D(H,\nu^\circ)$ and $\eta\in D(K,\nu)$, and let $y_1,y_2\in\mathcal U_0$. We then have
\begin{align*}
U_{H,K}^{\nu,\nu'}(\xi y_1\otimes_\nu y_2\eta)
&=U_{H,K}^{\nu,\nu'}\,(R^{\nu^\circ}(\xi)\otimes_\nu R^\nu(\eta))\,
\bigl(J\Lambda_\nu(y_1^*)\otimes_\nu\Lambda_\nu(y_2)\bigr)\\
&=(R^{\nu^\circ}(\xi)\otimes_{\nu'}R^\nu(\eta))\,
[U_{L^2(N)}^{\nu'}]^*\sigma^\nu_{-i/2}(y_1)\Lambda_\nu(y_2).
\end{align*}
Thus $U_{H,K}^{\nu,\nu'}$ is unambiguously determined on vectors of the form $\xi y_1\otimes_\nu y_2\eta$ as above, which, by \hyperref[p:2.2]{Remark~2.2}, form a total set in $H\otimes_\nu K$.

(b) By \cite[Proposition~3(b)]{1}, every operator in $\Lop_N(K,K')$ is a limit, in the topology of pointwise convergence, of linear combinations of operators of the form $t'^*t$, where $t\in\Lop_N(K,L^2(N))$ and $t'\in\Lop_N(K',L^2(N))$. By the separate continuity of the tensor product (\hyperref[p:2.2]{Remark~2.2}), it suffices to check that \hyperlink{diagram:2.6b}{diagram (b)} commutes for $T_1=t_1'^*t_1$ and $T_2=t_2'^*t_2$, with
\begin{align*}
t_1&\in\Lop_{N^\circ}(H,L^2(N)),&
t_1'&\in\Lop_{N^\circ}(H',L^2(N)),\\
t_2&\in\Lop_N(K,L^2(N)),&
t_2'&\in\Lop_N(K',L^2(N)):
\end{align*}
we then apply \hyperlink{part:2.6a}{(a)}.
\end{proof}

The ``tensor product over $N$'' is therefore independent of the weight $\nu$, as one might expect. The construction is entirely determined by its functorial properties and the equation $L^2(N)\otimes_N L^2(N)=L^2(N)$, in the sense that any notion of a ``tensor product over $N$'' satisfying this property and that of \hyperref[p:2.3]{Lemma~2.3} would necessarily be this one.

We give below an explicit construction of the isomorphism $U_{H,K}^{\nu,\nu'}$ in a particular case:

\itemhead{2.7}{Proposition.}
With the notation of \hyperref[p:2.6]{Proposition~2.6}, if $\nu'$ is the weight $\nu(h\,\cdot)$, where $h$ is a positive operator affiliated with the centralizer $N_\nu$ of the weight $\nu$, then the isomorphism $U_{H,K}^{\nu,\nu'}$ is given by
\begin{align*}
U_{H,K}^{\nu,\nu'}(\xi\otimes_\nu\eta)
&=\xi\otimes_{\nu'}h^{1/2}\eta,\\
&\hspace{-1em}\forall\xi\in H,\quad\forall\eta\in D(K,\nu)\cap\mathcal D(h^{1/2}),
\end{align*}
or equivalently,
\begin{align*}
U_{H,K}^{\nu,\nu'}(\xi\otimes_\nu\eta)
&=\xi h^{1/2}\otimes_{\nu'}\eta,\\
&\hspace{-1em}\forall\xi\in D(H,\nu^\circ)\cap\mathcal D(h^{1/2}),\quad\forall\eta\in K.
\end{align*}

\begin{proof}
Notice that, if $e$ is a spectral projection of $h$ corresponding to a bounded interval, and if $\eta$ is $\nu$-bounded, then $e\eta$ is $\nu$-bounded and lies in the domain of $h^{1/2}$: the stated properties therefore do characterize $U_{H,K}^{\nu,\nu'}$. Moreover, they make sense because, if $\eta\in\mathcal D(h^{1/2})\cap D(K,\nu)$, then $h^{1/2}\eta$ is $\nu'$-bounded and we have
\[
R^{\nu'}(h^{1/2}\eta)=R^\nu(\eta).
\]
It follows that
\[
\ip{\xi_1\otimes_\nu\eta_1}{\xi_2\otimes_\nu\eta_2}
=\ip{\xi_1\otimes_{\nu'}h^{1/2}\eta_1}{\xi_2\otimes_{\nu'}h^{1/2}\eta_2},
\]
for all $\xi_1,\xi_2\in H$ and all $\eta_1,\eta_2\in D(K,\nu)\cap\mathcal D(h^{1/2})$, and that there exists an isomorphism $W$ from $H\otimes_\nu K$ onto $H\otimes_{\nu'}K$ sending the tensor $\xi\otimes_\nu\eta$ to $\xi\otimes_{\nu'}h^{1/2}\eta$, for $\xi\in H$ and $\eta\in D(K,\nu)\cap\mathcal D(h^{1/2})$.

Let $W'$ be the isomorphism from $H'\otimes_\nu K'$ onto $H'\otimes_{\nu'}K'$ defined in the same way for a right $N$-module $H'$ and a left $N$-module $K'$. Then the diagram\tfnote{In the original, the upper-right object in this diagram is printed as $H'\otimes_{\nu'}K'$. It has been corrected to $H\otimes_{\nu'}K$, in accordance with the definition of $W$ and the vertical arrow.}
\[
\begin{tikzcd}[column sep=huge,row sep=large]
H\otimes_\nu K\arrow[r,"W"]\arrow[d,"T_1\otimes_\nu T_2"']
&H\otimes_{\nu'}K\arrow[d,"T_1\otimes_{\nu'}T_2"]\\
H'\otimes_\nu K'\arrow[r,"W'"']&H'\otimes_{\nu'}K'
\end{tikzcd}
\]
commutes for $T_1\in\Lop_{N^\circ}(H,H')$ and $T_2\in\Lop_N(K,K')$.

By the \hyperref[p:2.6]{preceding proposition}, it suffices to verify that $W=U_{H,K}^{\nu,\nu'}$ in the single case where $H$ is the right $N$-module $L^2(N)$ and $K$ the left $N$-module $L^2(N)$, a verification that presents no difficulty.
\end{proof}


\section{Standard Form of a von Neumann Algebra}\label{sec:3}

Let $\HH$ be a Hilbert space and $M\subset\Lop(\HH)$ a von Neumann algebra. We consider the following objects:
\begin{itemize}
\item the commutant $M'$ of $M$ in $\Lop(\HH)$;
\item the von Neumann algebra $N$ opposite to $M'$ (so that $\HH$ is an $M$-$N$ bimodule);
\item a normal semifinite faithful weight $\nu$ on $N$ and the corresponding weight $\mu'$ on $M'$;
\item the $M$-$M$ bimodule $\HH\otimes_\nu\overline{\HH}$;
\item the antilinear isometric involution $J_\nu$ of $\HH\otimes_\nu\overline{\HH}$ characterized by\tfnote{In this list, the original writes $D(\HH,\nu)$ twice. Both occurrences have been corrected to $D(\HH,\nu^\circ)$, following the convention for right $N$-modules in \hyperref[p:0.4.2]{0.4.2}. In \hyperlink{proof:3.1case4}{part~4) of the proof below}, the ordinary $H$ denoting this same Hilbert space has been normalized to $\HH$.}
\[
J_\nu(\xi\otimes_\nu\bar\eta)=\eta\otimes_\nu\bar\xi,
\qquad \forall\xi,\eta\in D(\HH,\nu^\circ);
\]
\item the closed cone $[\HH\otimes_\nu\overline{\HH}]_+$ generated by vectors of the form $\xi\otimes_\nu\bar\xi$, $\xi\in D(\HH,\nu^\circ)$.
\end{itemize}

\itemhead{3.1}{Proposition (Generalized Hilbert--Schmidt theorem).}
\emph{The quadruple $(\HH\otimes_\nu\overline{\HH},M,J_\nu,[\HH\otimes_\nu\overline{\HH}]_+)$ is a standard form of $M$ in the sense of \cite[Definition~2.1]{4}.}

\begin{proof}
\hypertarget{proof:3.1case1}{}1) Suppose that $M=N$, $\HH=L^2(N)$: identifying $\HH$ with $\overline{\HH}$ by the canonical involution $J$, it is easy to check that the isomorphism $U^\nu_{L^2(N)}$ from $L^2(N)\otimes_\nu L^2(N)$ onto $L^2(N)$ carries the involution $J_\nu$ to $J$ and the cone $[L^2(N)\otimes_\nu L^2(N)]_+$ to $L^2(N)_+$.

2) Suppose that $M$ is purely infinite: there exists an isomorphism of right $N$-modules from $\HH$ onto $q[L^2(N)]$, where $q$ is a projection in $N$. Thus $M$ is isomorphic to the reduced algebra $N_q$, and the proposition follows from \hyperlink{proof:3.1case1}{case~1)} by \cite[Lemma~2.6]{4}.

3) We are now reduced to the case where $M$ is semifinite. We can further reduce the problem by checking that the isomorphism $U^{\nu,\nu'}_{H,K}$ made explicit in \hyperref[p:2.7]{Proposition~2.7} carries the involution $J_\nu$ to $J_{\nu'}$ and the cone $[\HH\otimes_\nu\overline{\HH}]_+$ to $[\HH\otimes_{\nu'}\overline{\HH}]_+$. We may therefore restrict ourselves to the case where $\nu$ is a trace.

\hypertarget{proof:3.1case4}{}4) Let $\nu$ be a normal semifinite faithful trace on $N$. Let $\mu'$ be the corresponding trace on $M'=N^\circ$. There exists a normal semifinite faithful trace $\mu$ on $M$ such that $\frac{d\mu'}{d\mu}=1$ in the sense of \cite{1}. Applying the results of \hyperref[sec:1]{\S\,1}, we compute, for $\xi,\eta,\xi',\eta'\in D(\HH,\mu)=D(\HH,\nu^\circ)$:
\begin{align*}
\bigl\langle\xi\otimes_\nu\bar\eta,\xi'\otimes_\nu\bar\eta'\bigr\rangle
&=\bigl\langle\xi\langle\bar\eta,\bar\eta'\rangle_\nu,\xi'\bigr\rangle\\
&=\bigl\langle\langle\eta',\eta\rangle_{\mu'}\xi,\xi'\bigr\rangle\\
&=\bigl\langle\theta^\mu(\eta',\eta)\xi,\xi'\bigr\rangle\\
&=\bigl\langle R^\mu(\eta)^*\xi,R^\mu(\eta')^*\xi'\bigr\rangle
\qquad\text{(by \hyperref[p:1.7]{1.7})};
\end{align*}
hence there exists an isomorphism $W$ from $\HH\otimes_\nu\overline{\HH}$ onto $L^2(M)$ which assigns to a tensor $\xi\otimes_\nu\bar\eta$ as above the vector $R^\mu(\eta)^*\xi$ of $L^2(M)$. Now, by \hyperlink{eq:1.5b}{Lemma~1.5(b)} and \hyperref[p:1.2]{1.2}, we have
\[
R^\mu(\eta)^*\xi=\frac{d\omega_{\xi,\eta}}{d\mu}.
\]
Thus $W$ is clearly an isomorphism of $M$-$M$ bimodules\tfnote{The original prints ``$M$-$N$ bimodules.'' This has been corrected to ``$M$-$M$ bimodules,'' the structures carried by both $\HH\otimes_\nu\overline{\HH}$ and $L^2(M)$ in this statement.} which carries the involution $J_\nu$ to $J$ and the cone $[\HH\otimes_\nu\overline{\HH}]_+$ to $L^2(M)_+$.
\end{proof}

\itemhead{3.2}{Remark.}
Now that we know a priori that $\HH\otimes_\nu\overline{\HH}$ is standard for $M$, and in particular that the cone $[\HH\otimes_\nu\overline{\HH}]_+$ is self-dual, it is easy to check that the standard-form isomorphism $W$ from $\HH\otimes_\nu\overline{\HH}$ onto $L^2(M)$ is characterized as follows: for every normal semifinite faithful weight $\varphi$ on $M$, every $\xi$ in $\HH$, and every
\[
\eta\in D(\HH,\nu^\circ)\cap
\mathcal D\left(\left(\frac{d\varphi}{d\nu^\circ}\right)^{1/2}\right),
\]
we have
\[
W[\xi\otimes_\nu\bar\eta]
=R^\varphi\left(\left(\frac{d\varphi}{d\nu^\circ}\right)^{1/2}\eta\right)^*\xi.
\]

We thus obtain a remarkable analogy with the case of discrete factors, and a generalization of the results of \cite{8} for discrete algebras. Nevertheless, however satisfactory it may be from a formal standpoint, our \hyperref[p:3.1]{Proposition~3.1} should not be misleading: all the difficulty inherent in identifying standard forms is implicitly contained in the definition of the relative tensor product.

As usual, amplification and reduction yield the general commutation theorem from the standard case:

\itemhead{3.3}{Proposition.}
\emph{Let $H$ be a right $N$-module and $K$ a left $N$-module. In the Hilbert space $H\otimes_\nu K$, the commutant of $\Lop_{N^\circ}(H)\otimes_\nu\mathbb C_K$ is equal to $\mathbb C_H\otimes_\nu\Lop_N(K)$.}

As a second corollary of \hyperref[p:3.1]{Proposition~3.1}, we obtain that the procedure of relative amplification exhausts the representation theory of a von Neumann algebra:

\itemhead{3.4}{Corollary.}
\emph{Let $\HH$ be a Hilbert space and $M\subset\Lop(\HH)$ a von Neumann algebra of operators on $\HH$; denote by $N$ the algebra opposite to the commutant of $M$ in $\Lop(\HH)$, and fix a normal semifinite faithful weight $\nu$ on $N$.}

\emph{If $H$ is a left $M$-module, there exists an $N$-$\Lop_M(H)^\circ$ bimodule $K$, unique up to bimodule isomorphism, such that the $M$-$\Lop_M(H)^\circ$ bimodules $H$ and $\HH\otimes_\nu K$ are isomorphic.}

\begin{proof}
\emph{Existence.} Take $K=\overline{\HH}\otimes_\varphi H$, where $\varphi$ is a normal semifinite faithful weight on $M$.

\emph{Uniqueness.} If the $M$-$\Lop_M(H)^\circ$ bimodules $H$ and $\HH\otimes_\nu K$ are isomorphic, the $N$-$\Lop_M(H)^\circ$ bimodules\tfnote{In the first tensor product in the following display, the original prints $\nu$. This has been corrected to $\varphi$, the weight on $M$ used in the existence argument and in the next tensor product.}
\[
\overline{\HH}\otimes_\varphi H
\quad\text{and}\quad
\overline{\HH}\otimes_\varphi(\HH\otimes_\nu K)
=L^2(N)\otimes_\nu K=K
\]
are isomorphic.
\end{proof}

\itemhead{3.5}{Remark (Correspondences and functors).}
The results obtained allow us to reformulate, in terms of ``Hilbert-space bimodules,'' M.~A.~Rieffel's results on functorial correspondences between von Neumann algebras:

1) Let $M$ and $N$ be two von Neumann algebras. It is known (\cite[Proposition~5.4]{7}) that a functor $F$ (additive, normal, self-adjoint, $\ldots$) from the category $\mathit{normod}\,N$ of normal representations of $N$ to $\mathit{normod}\,M$ is characterized by the $M$-$N$ bimodule $F(L^2(N))$. The relative tensor product gives an explicit procedure for reconstructing the functor from the bimodule:

\itemhead{3.5.1}{Proposition (cf. \cite[Proposition~5.4]{7}).}
\emph{Every functor $F$ from $\mathit{normod}\,N$ to $\mathit{normod}\,M$ is equivalent to a functor of the form}
\begin{align*}
\{N\text{-modules}\}\ni K&\longmapsto H\otimes_\nu K,\\
L_N(K_1,K_2)\ni T&\longmapsto 1_H\otimes_N T,
\end{align*}
\emph{where $H$ is the $M$-$N$ bimodule $F(L^2(N))$. ($\nu$ is an arbitrary normal semifinite faithful weight on $N$.)}

2) It is easy to check that, for a given $M$-$N$ bimodule $H$, the functor of \hyperref[p:3.5.1]{Proposition~3.5.1} is an equivalence if and only if, on the one hand, $M$ and $N$ are faithfully represented and, on the other hand, the image of $N^\circ$ in $L(H)$ is equal to the commutant of the image of $M$ (cf.\ \cite[Theorem~8.15]{7}).

3) Functors and completely positive maps.

It is easy to make the connection with the corresponding examples introduced in \cite[\S\,I]{2} (and also to check that composition of functors is given by the ``relative'' tensor product of bimodules, that is, by composition of correspondences).

For example, if $\rho$ is a normal $*$-homomorphism from $M$ to $N$ (then $\rho(1)$ is a projection in $N$), the functor corresponding to the $M$-$N$ bimodule $L^2(\rho)$ is the one that, for a left $N$-module $K$, consists in reducing $K$ by the projection $\rho(1)$ and representing $M$, via $\rho$, on the space $\rho(1)\cdot K$.

More generally, if $P$ is a completely positive map from $M$ to $N$, and if $\HH$ is the bimodule associated with $P$ in \cite[Proposition~6]{2}, the functor $F$ from $\operatorname{Rep}N$ to $\operatorname{Rep}M$ associated with the $M$-$N$ bimodule\tfnote{The original prints ``$N$-$M$ bimodule.'' This has been corrected to ``$M$-$N$ bimodule,'' consistently with \hyperref[p:3.5.1]{Proposition~3.5.1} and the stated direction of the functor.} $\HH$ is what we called in \cite[\S\,2.3]{8} the Stinespring functor associated with $P$ (with reference to \cite{9}).

Thus every cyclic bimodule induces a Stinespring functor. More precisely, if $H$ is an $M$-$N$ bimodule, $\Omega$ a cyclic vector (i.e., $M\Omega N$ is total in $H$), and $\nu$ a normal semifinite faithful weight on $N$ such that $\Omega$ is $\nu^\circ$-bounded, then the functor associated with $H$ in \hyperref[p:3.5.1]{3.5.1} is equivalent to the Stinespring functor associated with the completely positive map from $M$ to $N$:
\[
M\ni x\longmapsto\langle x\Omega,\Omega\rangle_{\nu^\circ}\in N.
\]

As a corollary, every functor from $\mathit{normod}\,N$ to $\mathit{normod}\,M$ is, up to equivalence, a direct sum of Stinespring functors.

\begin{thebibliography}{9}
\phantomsection\addcontentsline{toc}{section}{References}\label{sec:references}
\bibitem{1} \textsc{Connes, A.}, On the spatial theory of von Neumann algebras, \emph{J. Functional Analysis}, \textbf{35} (1980), 153--164.
\bibitem{2} \textsc{Connes, A.}, Notes manuscrites, 1980.
\bibitem{3} \textsc{Dixmier, J.}, \emph{Les alg\`ebres d'op\'erateurs dans l'espace hilbertien}, Gauthier-Villars, Paris, 1969.
\bibitem{4} \textsc{Haagerup, U.}, The standard form of von Neumann algebras, \emph{Math. Scand.}, \textbf{37} (1975), 271--283.
\bibitem{5} \textsc{Hilsum, M.}, Les espaces $L^p$ d'une alg\`ebre de von Neumann, \emph{J. Functional Analysis}, \textbf{40} (1981), 151--169.
\bibitem{6} \textsc{Paschke, W. L.}, Inner product modules over $B^*$-algebras, \emph{Trans. Amer. Math. Soc.}, \textbf{182} (1973), 443--468.
\bibitem{7} \textsc{Rieffel, M. A.}, Morita equivalence for $C^*$-algebras and $W^*$-algebras, \emph{J. Pure Appl. Algebra}, \textbf{5} (1974), 51--96.
\bibitem{8} \textsc{Sauvageot, J. L.}, Produits tensoriels de $Z$-modules, \emph{Publ. Univ. P. \& M. Curie}, 23/1980.
\bibitem{9} \textsc{Stinespring, W. F.}, \href{https://doi.org/10.1090/S0002-9939-1955-0069403-4}{Positive functions on $C^*$-algebras}, \emph{Proc. Amer. Math. Soc.}, \textbf{6} (1955), 211--216.\tfnote{The original gives the year as 1975. This has been corrected to 1955, as recorded by the American Mathematical Society. The article title links to its official DOI.}
\end{thebibliography}

\begin{flushright}
\textit{JEAN-LUC SAUVAGEOT}\\
\textit{Universit\'e Paris VI,}\\
\textit{Lab. de Calcul des Probabilit\'es,}\\
\textit{Tour 56--4, Place Jussieu,}\\
\textit{75230 Paris Cedex 05,}\\
\textit{France.}
\end{flushright}

\begin{center}
Received June 20, 1981; revised June 15, 1982.
\end{center}

\end{document}
